\chapter{Device Endpoint Vulnerabilities}
\label{ch:4.1}
\chaptertopic{Topic 4.1}

\chapterquote{Only amateurs attack machines; professionals target people.}{Bruce Schneier, ``Semantic Attacks: The Third Wave of Network Attacks'' (2000)}

\section{Motivation}

Units 2 and 3 protected space and wire. Unit~4 moves to the endpoint itself -- the laptop, phone,
camera, or thermostat where data actually lives and computation actually happens. This chapter
covers three distinct ways an endpoint gets compromised: malicious software running on it, and
physical tampering with the hardware itself. It also covers the enormous, largely undefended
category of Internet of Things devices, which fail in both of the first two ways simultaneously,
at a scale no other device category matches.

\begin{apalignment}{Unit 4: Securing Devices}{Topic 4.1 -- Device Vulnerabilities and Attacks}
\begin{apobjectives}
\apobjective{4.1.A}{Analyze Risk}{Recognize the different classes of computing device -- servers,
workstations and laptops, mobile handsets, network appliances, and embedded/industrial controllers
with no screen or keyboard -- and account for how each class differs in what it can run, what it
can be patched with, and how long its vendor will support it.}
\apobjective{4.1.B}{Analyze Risk}{Work out, from what an incident actually did, which family of
malicious software was involved: whether the code needed a host file and a human, whether it moved
between machines unaided, whether it arrived disguised as legitimate software, and whether its
purpose was extortion, concealment, surveillance, or theft.}
\apobjective{4.1.C}{Analyze Risk}{Trace the route from a weakness in a device -- a factory
password, an unneeded service, an unpatched kernel, unsupported firmware -- to the concrete harm
an adversary can produce with it: data lost, hardware damaged, operations halted, records
destroyed.}
\apobjective{4.1.D}{Analyze Risk}{Evaluate the weaknesses present on a fleet of devices and record
them in a form someone else can act on -- each finding scored for likelihood and impact, ranked
against the others, and documented where the decision to fix or accept it can later be audited.}
\end{apobjectives}
\end{apalignment}

\section{Device Classes and the Support Lifecycle}

Not every endpoint is the same kind of target. Five device classes recur across a typical
organization, and each one differs in what it can run, how it gets patched, and how long its
vendor keeps patching it at all.

\begin{itemize}
\item \textbf{Servers} run general-purpose operating systems and patch on a scheduled,
change-controlled basis -- downtime costs money, so a patch waits for a maintenance window rather
than shipping the moment a fix exists.
\item \textbf{Workstations and laptops} run the same kind of general-purpose operating system but
patch far more aggressively, often automatically, because a compromised end-user machine is
usually the first hop into the rest of the network.
\item \textbf{Mobile handsets} run a locked-down operating system controlled by one vendor, which
pushes security updates directly to the device. That vendor also sets its own support cutoff,
though: past that date, even a healthy phone stops receiving patches entirely.
\item \textbf{Network appliances} -- routers, firewalls, switches -- run stripped-down firmware
built for one function. An administrator updates it rarely, and only by scheduling an outage.
\item \textbf{Embedded and industrial controllers} have no screen, no keyboard, and often no
practical way to patch at all. The firmware that shipped with the device may be the only firmware
it ever runs, across a service life measured in decades rather than years.
\end{itemize}

The pattern across all five classes is the same: the harder a device is to patch, and the longer
it stays in service, the longer any given vulnerability in it stays exploitable. An unpatched
laptop is a bad week. An unpatched industrial controller can be a bad decade.

\section{Malware Types}

\begin{definition}{Malware}
Malicious software: any program intentionally designed to cause harm, steal data, disrupt
operations, or give an attacker unauthorized control.
\end{definition}

Four categories account for the overwhelming majority of real-world device compromises. Each has a
genuinely different mechanism, which is exactly why classifying an infection correctly matters --
the right defense is specific to the category, not generic.

A \textbf{virus} attaches its code to a legitimate host file -- an executable, a macro-enabled
document, a script -- and activates only when a user runs that host file. At that point it
typically searches for other files to infect, and spreads further whenever an infected file is
copied or emailed elsewhere. Its defining trait is dependence on a host file \emph{and} direct user
action to propagate. A \textbf{worm} shares the self-replicating instinct but needs neither: it is
a self-contained program that actively probes a network and finds vulnerable devices. It copies
itself across automatically, repeating the process on each new host, with no user interaction
required anywhere in the chain. The 2003 SQL Slammer worm shows how fast that unaided spread can
go: it doubled its infected-host count every 8.5 seconds by exploiting a single buffer-overflow
flaw in Microsoft SQL Server. Within its first ten minutes it reached an estimated 75,000
vulnerable hosts.\srccite{4.1}{1}

A \textbf{trojan} takes the opposite approach entirely: it does not self-replicate at all, relying
instead on disguise -- presenting itself as a free game, a pirated application, a ``system
cleaner'' -- so the user installs it willingly. Once installed, a trojan often opens a back door
granting an attacker remote control, keystroke logging, or a launch point for further attacks. And
because the user installed it themselves, tools that look for unauthorized changes may never flag
the initial infection at all. \textbf{Ransomware} encrypts a victim's files with a key only the
attacker holds, then demands payment, typically in cryptocurrency, for the decryption key. It
disproportionately targets hospitals, schools, and municipalities, precisely because they cannot
afford the downtime that refusing to pay would cost them.

\begin{figure}[h]
\centering
\begin{forest}
for tree={
  draw=kickstartermaroon, thick, rounded corners, align=center,
  font=\scriptsize, l sep=14pt, s sep=10pt, minimum height=0.7cm,
}
[Malware
  [Virus\\{\tiny needs host file + user}, fill=kickstartercream]
  [Worm\\{\tiny self-contained, no user}, fill=kickstartercream]
  [Trojan\\{\tiny disguise, no self-replication}, fill=kickstartercream
    [Ransomware\\{\tiny encrypts, demands payment}, fill=kickstartermaroon!15]
  ]
]
\end{forest}
\caption{The four malware categories divide on two independent questions: does it need a host
file and user action to spread, and does it rely on disguise rather than self-replication?
Ransomware is a payload, not a spreading mechanism -- it commonly arrives via a trojan.}
\label{fig:4.1-malware-tree}
\end{figure}

\begin{lstlisting}[language=Python, caption={Classifying an infection by its behavioral signature}]
SIGNATURES = {
    "unusual_network_scanning": "worm",
    "self_installed_via_deception": "trojan",
    "host_file_infection": "virus",
    "files_encrypted_ransom_note": "ransomware",
}

def classify_malware(signature):
    """Returns the malware category matching a behavioral signature
    string, or "unknown" if unrecognized."""
    return SIGNATURES.get(signature, "unknown")
\end{lstlisting}

These categories are not mutually exclusive in a real attack -- a worm might carry a trojan
payload that later downloads ransomware, stacking three mechanisms into one incident. Malware
enters through an \textbf{infection vector}: a phishing attachment (Chapter~\ref{ch:1.1}), a
drive-by download, an infected USB drive, or exploitation of an unpatched vulnerability. Because
the defense is vector-specific, recognizing which vector was used is not academic. No
attachment-scanning policy stops worm spread; no network segmentation stops a user from installing
a disguised trojan; only an offline backup neutralizes ransomware's actual leverage, since it
removes the ``pay or lose everything'' choice entirely.

\begin{misconception}
\textbf{``Ransomware destroys your files.''} Ransomware typically encrypts files, leaving them
fully intact but inaccessible without the key -- the attacker can decrypt them if paid, which is
exactly the leverage a backup breaks. The files are locked, not gone.
\end{misconception}

\section{Device Tampering}

The most fundamental rule of hardware security is blunt: if an attacker has physical access to
your device, it is no longer your device. Passwords, firewalls, and network access lists all
assume the physical chassis itself is secure. Once that assumption fails, almost every
software-based defense can be bypassed without ever touching the operating system's login prompt.

A \textbf{USB Rubber Ducky} exploits the fact that operating systems automatically trust a
keyboard the instant one is plugged in. The device is programmed to mimic a USB keyboard rather
than storage, and once connected it types scripted commands at over 1{,}000 words per
minute\srccite{4.1}{2} -- opening a terminal, downloading malware, exfiltrating data. None of it
looks like a file to scan, so antivirus never sees it. A \textbf{hardware keylogger} sits inline
between the keyboard cable and the computer, intercepting every electrical signal before it ever
reaches the operating system. That is precisely why no software-based scanner can detect it: the
interception happens entirely in hardware, one layer below anything software can observe.

\begin{example}[The cold boot attack]
When a computer powers off, its RAM does not lose its contents instantly -- it degrades over
several seconds to minutes. An attacker who sprays the memory chips with compressed-air duster held
upside down -- so it discharges a cold liquid instead of a gas -- can chill them to roughly $-50^\circ$C,
freezing that decay in place. At that temperature, research on the technique found that cutting
power for a full minute still let an attacker recover at least 99.9\% of the RAM's
bits.\srccite{4.1}{3} The attacker then quickly transfers the physical RAM module to another
machine before it fades, extracting cryptographic keys or passwords that were resident in memory
moments before. The attack works entirely on physics the operating system has no way to prevent:
by the time it runs, the OS that would have protected that memory is no longer even executing.
\end{example}

Physical attacks succeed because operating systems are built for convenience first. They trust a
newly connected keyboard implicitly, with no administrator privilege required. By default they
will also boot from whatever storage device is connected -- letting a live Linux USB override the
installed OS and every protection running inside it. \textbf{Full-disk encryption (FDE)} defeats
this specific path. Bound to a \textbf{Trusted Platform Module (TPM)} that verifies boot integrity
before releasing the decryption key, FDE ensures that an attacker who removes the drive or boots
from external media still cannot read the encrypted contents without the key. BIOS/UEFI password
protection and disabling external-media boot close the boot-integrity gap directly, while port
blockers, cable locks, and chassis intrusion switches raise the physical bar for opportunistic
tampering.

\begin{lstlisting}[language=Python, caption={Detecting physical tampering events in a hardware audit log}]
TAMPER_EVENTS = {"chassis_open", "unauthorized_usb_keyboard", "unknown_boot_device"}

def flag_tamper_events(log):
    """log: list of dicts like {"event": str, "timestamp": int}.
       Returns sorted timestamps of recognized tampering events."""
    return sorted(e["timestamp"] for e in log if e["event"] in TAMPER_EVENTS)
\end{lstlisting}

\begin{misconception}
\textbf{``Operating system passwords protect my data if my laptop is stolen.''} A login password
only guards the intended boot path. An attacker can remove the drive and read it on another
machine, or boot a Linux USB and browse the filesystem directly, unless the disk's contents are
themselves encrypted through FDE.
\end{misconception}

\section{IoT Device Vulnerabilities}

The Internet of Things -- cameras, routers, thermostats, door locks, medical sensors, industrial
controllers -- numbers in the tens of billions of deployed devices, and collectively they are the
softest target on the modern network. The reasons recur across nearly every device category, which
is what makes them worth understanding as a pattern rather than a checklist.

Most IoT devices ship with \textbf{default credentials} -- \texttt{admin}/\texttt{admin},
\texttt{root}/\texttt{12345}, or nothing at all -- meant only for initial setup. Most owners never
change them, because a smart light bulb does not present as ``a computer that needs a password.''
Automated scanners try these well-known lists against every IP address on the internet and succeed
within seconds wherever the default was never touched. \textbf{Firmware} -- the low-level software
controlling a device's core operation -- typically has no automatic update mechanism. Many
manufacturers stop supporting a device (it becomes \textbf{end-of-life}) within a year or two. Any
undiscovered vulnerability then stays exploitable for the rest of the device's operational life --
often five to ten years for a security camera or industrial sensor.

Devices frequently expose \textbf{management interfaces} over old, insecure protocols to keep
setup convenient. Telnet, a plaintext remote-login protocol, is one; UPnP is another, letting a
device silently punch a hole through a home router's firewall without the owner's knowledge and
exposing an administrative interface directly to the internet. Their companion mobile apps depend
on \textbf{cloud APIs} that, if improperly authenticated, let an attacker control or query devices
belonging to other users entirely. Traffic sent without TLS, meanwhile, is readable by anyone
sharing the same network segment. Physical accessibility compounds all of this: a doorbell camera
on an exterior wall or a sensor in a parking garage is reachable by anyone who can touch it. That
often lets an attacker extract firmware directly from a chip or attach a debug interface,
bypassing network defenses entirely.

\begin{example}[Mirai and the aggregate risk of scale]
Any single insecure smart plug looks low-stakes in isolation. The 2016 Mirai botnet showed why
scale changes that calculus completely. It scanned the internet for IoT devices still using
factory-default Telnet credentials and logged in automatically wherever it found one. That
enlisted each device into a network that reached over 600{,}000 compromised devices under central
control at its peak.\srccite{4.1}{4} Mirai used that aggregate to launch some of the largest
distributed-denial-of-service attacks ever recorded, including the October 2016 attack on DNS
provider Dyn that briefly knocked Twitter, Netflix, and other major sites offline. No individual
device mattered; the aggregate was devastating. That is exactly why security teams increasingly
worry about \textbf{shadow IoT} -- devices connected to a network that IT never approved or
secured, silently expanding the attack surface with nobody tracking it.
\end{example}

\begin{lstlisting}[language=Python, caption={Auditing an IoT device's primary security weakness}]
def audit_device(device):
    """device: dict like {"default_credentials": bool, "firmware_current": bool,
       "telnet_open": bool, "upnp_enabled": bool, "api_authenticated": bool,
       "traffic_encrypted": bool}.
       Returns the highest-priority weakness label, in priority order."""
    if device["default_credentials"]:
        return "default-credentials"
    if not device["firmware_current"]:
        return "outdated-firmware"
    if device["telnet_open"] or device["upnp_enabled"]:
        return "exposed-management-interface"
    if not device["api_authenticated"]:
        return "unauthenticated-api"
    if not device["traffic_encrypted"]:
        return "unencrypted-traffic"
    return "no-major-weakness"
\end{lstlisting}

\begin{misconception}
\textbf{``If a device isn't connected to the internet, it's safe.''} A device on a local network
can still be reached by anything else already on that network -- a compromised laptop, a malicious
guest device -- and many ``local-only'' devices still phone home to a cloud API regardless. Network
placement reduces risk; it does not eliminate it.
\end{misconception}

\section*{Review Questions}
\begin{reviewquestions}
\item A hospital patches its laptop fleet automatically within days of every vendor security
bulletin, but its infusion-pump controllers are still running the firmware they shipped with a
decade ago. Explain why these two device classes end up with such different exposure windows,
referencing what each class can run and how (or whether) it gets patched.
\item A university network shows thousands of new infections within minutes, with no phishing
email or user download involved in any of them. Identify which malware category this matches and
explain the specific mechanical property that makes this scale of spread possible without user
action.
\item Explain why antivirus software cannot detect a hardware keylogger, tracing the explanation to
exactly where in the system the interception occurs relative to the operating system.
\item A stolen laptop has a strong login password but no full-disk encryption. Explain precisely
how an attacker extracts its data anyway, and identify the single control that would have
prevented it.
\item An IoT device has default credentials, outdated firmware, and unencrypted traffic
simultaneously. Using the priority-ordered audit logic from this chapter, explain which weakness
should be reported first and why the ordering itself is a meaningful design decision, not arbitrary.
\item Explain why the Mirai botnet is described as an aggregate risk rather than an individual-device
risk, and explain what specific IoT weakness it exploited to reach its scale.
\item Using the malware taxonomy diagram, explain why ``ransomware'' and ``trojan'' are not
mutually exclusive categories the way ``virus'' and ``worm'' are, and describe a realistic
infection chain that involves both labels correctly at once.
\item A security team has a strong attachment-scanning policy and a strong network-segmentation
policy but no offline backups. Using the vector-specific reasoning from this chapter, explain
which malware category this fleet remains fully exposed to and why neither existing control
addresses it.
\end{reviewquestions}

\begin{answerkey}
\item Workstations and laptops run general-purpose operating systems with vendor-managed, often
automatic patch delivery, so a newly disclosed vulnerability's exposure window is measured in
days. Embedded and industrial controllers such as infusion pumps run purpose-built firmware with
no screen or keyboard, often no practical update mechanism at all, and a service life measured in
decades rather than years -- so a vulnerability discovered once can stay exploitable for the
device's entire remaining life. The gap is not about which device matters more; it is about what
each device class can run, and whether that software can realistically be replaced after the
device is deployed.
\answerlabel{4.1.A}
\item The infection matches a \textbf{worm}. A worm is a self-contained program that actively
probes a network for vulnerable hosts and copies itself across automatically, with no host file and
no user action required anywhere in the chain -- which is exactly the mechanical property that lets
an infection count climb into the thousands within minutes, the same doubling behavior the 2003
Slammer worm showed.
\answerlabel{4.1.B}
\item A hardware keylogger sits inline on the physical cable between the keyboard and the
computer, intercepting the electrical keystroke signal before it ever reaches the operating system.
Antivirus software runs as code \emph{within} the operating system, one layer above where the
interception happens, so it has no vantage point from which the tampering is even visible.
\answerlabel{4.1.C}
\item A login password only protects the intended boot path. The attacker can simply remove the
drive and read it on another machine, or boot the laptop from a live Linux USB and browse the
filesystem directly. Either path bypasses the login prompt entirely, because the data on disk was
never encrypted. Full-disk encryption (FDE) is the single control that would have prevented this, since it
makes the drive's contents unreadable without the key regardless of which machine reads the physical
disk.
\answerlabel{4.1.C}
\item Default credentials should be reported first. The audit logic in this chapter checks it
first because it is both the easiest weakness to exploit (no skill or exploit code needed, only a
known username/password pair) and the one automated internet-wide scanners actively search for at
scale, whereas outdated firmware or unencrypted traffic both require more attacker effort to
leverage. The ordering reflects likelihood of exploitation, not an arbitrary list position.
\answerlabel{4.1.D}
\item Mirai is an aggregate risk because no single insecure smart plug or camera caused meaningful
harm on its own; the damage came from combining hundreds of thousands of such devices under one
attacker's control to launch massive distributed denial-of-service attacks. It exploited factory
default Telnet credentials that device owners never changed, letting Mirai's scanner log in
automatically at internet scale.
\answerlabel{4.1.B}
\item Virus and worm differ on a spreading mechanism -- host-file-and-user versus self-contained --
so an infection is mechanically one or the other, never both at once. Ransomware and trojan differ
on a different axis entirely: trojan describes \emph{how a program gets installed} (disguise
rather than self-replication), while ransomware describes \emph{what the payload does once
running} (encrypt files and demand payment). A realistic chain: a user downloads a disguised
``free software cracker'' (a trojan, by installation method) which, once run, encrypts the user's
files and displays a ransom note (ransomware, by payload behavior) -- both labels correctly
describe the same single infection because they answer different questions about it.
\answerlabel{4.1.B}
\item The fleet remains fully exposed to ransomware. Attachment scanning targets the phishing
infection vector and network segmentation limits lateral movement after a compromise, but neither
addresses ransomware's actual leverage, which is the absence of any way to recover encrypted files
without paying. A well-disguised trojan can still be willingly installed by a user despite
attachment scanning (if it arrives by a different vector, such as a fake software download), and
once ransomware executes locally it does not need to move laterally to cause damage -- only offline
backups remove the ``pay or lose everything'' choice that is ransomware's entire mechanism of
harm.
\answerlabel{4.1.B}
\end{answerkey}

\section*{Sources}
\begin{enumerate}[label=\lbrack\arabic*\rbrack, leftmargin=2.2em]
\item \srcitem{4.1}{1}David Moore, Vern Paxson, Stuart Staniford, et al., ``The Spread of the
Sapphire/Slammer Worm,'' CAIDA (Cooperative Association for Internet Data Analysis), 2003.
\item \srcitem{4.1}{2}Hak5, ``USB Rubber Ducky'' product specifications, shop.hak5.org (accessed
2026).
\item \srcitem{4.1}{3}J. Alex Halderman, Seth D. Schoen, Nadia Heninger, et al., ``Lest We
Remember: Cold-Boot Attacks on Encryption Keys,'' \emph{Communications of the ACM}, vol.\ 52,
no.\ 5 (May 2009).
\item \srcitem{4.1}{4}Elie Bursztein, ``Inside Mirai the Infamous IoT Botnet: A Retrospective
Analysis,'' elie.net, 2017; peak device count corroborated by contemporaneous reporting on the
October 2016 Dyn DNS attack.
\end{enumerate}
